%%%%%%%%%%%%%%%%%%%%%%%%%%%%%%%%%%%%%%%%%%%%%%%%%%%%%%%%%%%
%                                                         %
% Cette version repart de 03_Draft\data_current-short.tex %
%                                                         %
%%%%%%%%%%%%%%%%%%%%%%%%%%%%%%%%%%%%%%%%%%%%%%%%%%%%%%%%%%%

% NB: Je supprime toutes les (grosses) parties qui étaient en commentaire dans la version précédente pour qu'on y voit plus clair.

\subsection{Data}
Throughout the paper, we rely on several French administrative sources. We first use linked employer-employee datasets (Panel DADS and FH-DADS) to identify workers, their job (wage, hours, employer) and layoff. Panel DADS is a random sample of 1/12 of the universe of workers, and FH-DADS matches this panel with unemployment insurance administrative records. 
% , we observe firm's unique, time-invariant identifier (SIREN code), which allows us to retrieve firms' balance-sheet and income statement from the FICUS-FARE datasets. 
We also use the exhaustive version of the linked employer-employee data, DADS Postes, to estimate firm fixed effects. %in section \ref{sec:data:firm:premium}
Firms' balance-sheet and income statement are retrieved from the FICUS-FARE administrative dataset, built from the corporate tax returns and the social security declarations of the universe of firms in the non-financial corporate sector. All datasets are merged through the firm identifier (SIREN).\footnote{See Appendix Table \ref{apptab:datasets} for more details about data.}
%In the remainder of this section, we explain how we build our samples and how we measure our main outcomes of interest. Appendix Table \ref{apptab:datasets} provides more details about our datasets.


\subsection{Displaced workers}  
\label{sec:data:displaced}


%n our main sample, we follow workers laid-off for economic reasons (\textit{licenciés économiques}, LE hereafter) between 2002 and 2012. 
We identify workers laid-off for economic reasons (\textit{licenciés économiques}, LE hereafter) using FH-DADS data, which matches: \textit{i)} a panel of workers with data on wages, hours, and employer (panel DADS) and \textit{ii)} unemployment insurance data with details on unemployment spells (registration date, duration, benefits). Crucially, this dataset, which covers the period from 2002 to 2012, precisely identifies LE workers by documenting the reason for registration at the French employment agency (e.g., layoff for personal vs. economic reasons) as declared by the firm to the administration.

Economic layoffs require firms to justify economic difficulties, as detailed in Appendix \ref{appsec:LE_def}. Given the strict legal definition and conditions, firms are unlikely to misdeclare personal layoffs as economic ones, despite potential incentives (\cite{fraisse2015labor} document that personal dismissals are more often contested in court).\footnote{Supporting this, \cite{signoretto2019individual} find that economic dismissals are largely explained by firms' economic conditions, while personal layoffs are linked to management style.}

Focusing on a sample of LE workers has two main advantages. First, it minimizes the risk of including individuals who voluntarily left their firm before or during a mass layoff. This concern has been emphasized by \cite{seim_2019}, who quantifies the bias resulting from not being able to separate quits from layoffs (\textit{e.g.} focusing on workers leaving the firm within a time window of an important downsizing of the labor force). Moreover, because economic layoffs in our setting are legally defined and reported by firms to the administration, our measure reflects the circumstances of separation rather than inferring displacement indirectly from firm-level employment changes. In this respect, our approach is related to other recent efforts to move beyond canonical mass-layoff definitions, including \cite{seim_2019,fackler2017explaining,cederlof_2024} and \cite{bertheau_etal_2023}.\footnote{Importantly, workers laid off for economic reasons are eligible for the standard unemployment insurance scheme applicable to all involuntarily displaced workers in France, with no specific difference in benefit generosity or access to labor market programs relative to other forms of layoffs; we provide further institutional details in Appendix \ref{appsec:LE_def}}

% NB, Seim 2019 dit :
% Finally, by using administrative data on displacements, it is possible to assess quantitatively the bias that results from previous studies not being able to separate quits from layoffs.
%The final part of the paper addresses the potential bias that might plague previous studies of effects of job loss that use matched employeremployee data and thus are unable to distinguish quits from layoffs. I use the Swedish employer-employee records and define separations of matches as job displacements according to the literature. Estimates from regressions using conventional approaches are compared with the estimates obtained when using the PES data, where displacements are directly observed. I do not only assess the potential bias in previous studies, but also propose how to define separations as displacements in the absence of direct data on job displacements. First, separations in periods of mass separations are defined as displaced, following e.g. (Jacobson et al., 1993) and (Sullivan and von Wachter, 2009). The mass layoff estimates turn out to significantly understate the effects of displacement. Instead, defining separations in three-year time windows prior to plant closure as displacements leads to an overestimation of the effects. The arguably most accurate estimates are obtained by defining displacements as separations occurring during a one-year window preceding plant closure.

Second, it allows for a more representative set of employers compared to the literature on mass layoffs, which typically defines them as a 30\% or more decline in labor force \citep[e.g.][]{Jacobson1993,Gathmann2018,Gulyas2019,Lachowska2020,Schmieder2020} and imposes a minimum-firm-size restriction (often $\geq 50$ employees). We deliberately do not impose such a size threshold: rather than firm size in headcount terms, the relevant dimension for our analysis is firm wage-setting, which we summarise with the AKM firm fixed effect. Along this dimension, our LE sample is closer to the population firm-premium distribution than the typical mass-layoff sample, which selects on relatively large and high-premium employers (Appendix Figure \ref{fig:sample_FEdist}).
% This footnote is now in the main text in the "rob test" subsection.
%\footnote{As noted by a referee, in small firms, firm-level difficulties may still partly reflect the performance of individual workers, so that layoffs classified as economic may in practice be related to worker-specific factors. Figure \ref{fig:wage_decomposition_size10} shows that our main results are robust to restricting our sample to firms with more than 10 workers.}

To ensure our sample definition is not the sole driver of results, we build a second sample based on mass layoffs (MLO): plants losing at least 30\% of employees. This approach aligns with the literature and allows for comparison with prior studies. Appendix \ref{appsec:mlo} provides details on this sample.

\paragraph{Construction of the treatment (LE workers) and control groups.} 

The year of layoff, denoted $t^D$, is when the individual registers for unemployment after an economic layoff. We define a LE worker's previous job as the position with the end date closest to the contract end date recorded by the unemployment agency. We focus on individuals losing stable jobs, selecting workers aged 20-50 at layoff %, who worked over 1800 hours in the previous job in $t^D-1$ and $t^D-2$, 
and who are not employed in the same firm on June 30, $t^D+1$. In addition, we impose a three year tenure condition, so as to analyze a sample of workers with a strong labor market attachment: LE workers must be observed in the firm at least on $t^D$, $t^D-1$ and $t^D-2$.\footnote{The first LE that are included in our analysis are observed in 2005.} % The latest LE workers that are included in the analysis are observed in 2012, but for them we do not have any "post-layoff" observation.
An individual is included only when first registered as LE and meeting these criteria. Overall 20,897 -> \textcolor{red}{21,132} individuals, displaced from 14,038 -> \textcolor{red}{14,115} distinct firms, fit these criteria. \textcolor{orange}{[21,132 and 14,115 OK: Table 1, last two rows]}
 
We use propensity score matching to construct a control group for LE workers. Matching is performed on all non-displaced workers in FH-DADS who meet the criteria applied to the LE workers in terms of job tenure. %and hours worked. 
Matching occurs within year and sector-size cells, defined by the intersection of 10 industries (1-digit) and 4 quartiles of firm size (based on the size distribution at $t^D-1$ -> \textcolor{red}{$t^D-2$} among LE). We match workers using age, firm size at $t^D-2$ -> \textcolor{red}{$t^D-1$} \textcolor{orange}{[$t^D-1$ OK: Table 1 row ``\# employees at firm at $t^D-1$''; the quartiles of the size distribution at $t^D-2$ cannot be checked against a table]}, and hourly wage and hours worked in the main job in $t^D-2$ and $t^D-3$. Using pre-displacement characteristics measured sufficiently far in advance helps limit the risk that these variables are already affected in the period immediately preceding displacement. We estimate a probit to predict the propensity score and keep the nearest neighbor, with replacement. The LE panel (LE workers and their control group) is described in Section \ref{sec:stat_desc}.\footnote{Our control group consists of workers who are never displaced through either an LE or an MLO30 event over the sample period; we do not require controls to remain continuously employed at the same firm over the post-event window---in line with \citet{Schmieder2020}, \citet{Lachowska2020} and \citet{bertheau_etal_2023}, and departing from the convention of \citet{Jacobson1993}. Appendix~\ref{app:stayer_controls} re-estimates the main results on the MLO30 sample using the JLS-style stayer control group, as a comparability check with the existing literature.}

\subsection{Firm wage premia}
\label{sec:data:firm:premium} 

We compute firm wage premia using the exhaustive version of the matched employer-employee data (DADS \textit{postes}). This version of DADS allows us to estimate firm fixed effects for a larger set of firms compared to studies focusing on the worker-level panel \citep[e.g.][]{Abowd1999,bertheau_etal_2023} which is a 1/12th sample of the exhaustive file. In the exhaustive dataset, worker identifier is renewed at each vintage but any vintage $t$ contains information on jobs at time $t$ and $t-1$. Following \cite{godechot2023akm}, we take advantage of this overlap, using common individual information between vintage $t$ and $t-1$ in order to build a new worker panel based on the exhaustive data, which covers about 98\% of all workers between 2002 and 2015.

The DADS \textit{postes} presented here and the panel DADS cannot be merged at the worker-level. This implies that we cannot recover worker fixed-effects. We view this limitation as innocuous as we are primarily interested in the role of firms characteristics in explaining the cost of displacement---unlike for instance \citep{helm2022} who study how this effect varies with worker-fixed effect.

The AKM-type model we fit writes as follows:
\begin{equation}\label{eq:akm_basics}
	\ln \left(w_{it}\right)=\psi_{J(i, t)}+\alpha_{i}+\lambda_{t}+\epsilon_{i t}, 
\end{equation}
where $w_{it}$ refers to worker $i$'s  hourly wage in year $t$, \textcolor{red}{winsorized within year at 1\%/99\%} \textcolor{orange}{[OK: the AKM code winsorizes the hourly wage at 1/99 year by year]}, $\psi_{J(i, t)}$ is the fixed effects of firm $J$ where worker $i$ was employed at period $t$ and where $\alpha_{i}$ and $\lambda_{t}$ are individual and time fixed effect respectively. The period used to estimate firm fixed effects covers years 2002 to 2015.

To limit contamination of firm fixed effects by displacement-induced worker reallocation, we exclude from the AKM estimation sample all post-event observations of workers leaving an LE or MLO firing firm in the year of the event or the year after, while retaining their pre-event observations. Because the worker identifiers used in the event-study sample (\textit{panel DADS}) cannot be directly linked to those used in the exhaustive matched employer--employee data (\textit{DADS postes}), this exclusion is implemented at the firm-year level rather than for displaced workers in our estimation sample only. Concretely, any worker leaving an LE or MLO firm in year $t^D$ or $t^D+1$ is excluded from the AKM sample from that point onward. This procedure follows the logic of the contamination corrections, which typically exclude post-displacement observations from the AKM estimation sample.\footnote{Recent studies have emphasized that displacement-induced worker reallocation can contaminate AKM firm fixed effects if post-displacement observations are included in the estimation sample. \citet{Lachowska2020} exclude displaced workers entirely from the AKM estimation sample, while \citet{moore2019firm} and \citet{Schmieder2020} exclude post-layoff observations only. Our approach is closest in spirit to the latter studies, although the absence of a direct worker-level link between \textit{panel DADS} and \textit{DADS postes} requires implementing the exclusion at the firm-year level.} However, our exclusion rule is more stringent than in the existing literature, since it applies to all workers leaving LE or MLO firms rather than only to displaced workers included in the event-study sample.

Firm and worker fixed effects have been shown to suffer from limited mobility bias \citep{andrews2008high}, which implies an upward bias on the estimate of the variance of firm fixed effects and a downward bias regarding the covariance between worker and firm fixed effects. Our use of the near universe of across-firms movement tempers this risk. Moreover, this type of bias is particularly relevant when studying assortative matching \citep{bonhomme2019distributional} or the contribution of firm to wage dispersion \citep{song2019firming}. %{Recent contributions have proposed either alternative estimation framework to the AKM two-way fixed effect approach \citep{bonhomme2019distributional,bonhomme2020much} or corrections applied to the variance and covariance estimator \citep{kline2020leave}.} 
In this paper, however, we are interested in the first moment of these firm fixed effects. Accordingly, measurement error occurs in our dependent variable and is likely to affect the precision of our estimates but not to bias them in any systematic way.\footnote{Recent work on French data corroborates this assessment for long panels: both \citet{godechot2023akm} and \citet{milsom_hou_2024} find that the Kline-Saggio-S{\o}lvsten correction for limited mobility bias is small in magnitude when AKM is estimated over time spans similar to ours, leaving the first moment of the firm fixed effects essentially unaffected. \textcolor{orange}{[reference check, 2026-11-20: neither paper reports this. \citet{godechot2023akm} use a split-sample correction, not KSS, and find significant differences between corrected and uncorrected variance components; Milsom and Hou's French sample covers 2015--2019, their paper is titled ``Do Workers Sort to Firms or to Occupations?'', and no paper with the cited title exists. Proposed: ``Limited mobility bias affects quadratic functions of the estimated firm effects, such as their variance or their covariance with worker effects \citep{andrews2008high}, but the estimated effects themselves are unbiased \citep{godechot2023akm}. Since our outcome is the firm effect itself, the bias adds noise to the dependent variable without shifting the event-study coefficients.'' Then drop Milsom and Hou.]}}


\subsection{Other firm characteristics}
\label{sec:data:firm:other}

We measure labor productivity and labor share using the FICUS-FARE database, as processed by \citet{burstein_etal_2025} (see Appendix Table \ref{apptab:datasets}). We build two productivity measures: value added per worker, and a measure adjusted for capital intensity. The latter is the residual of a regression of value added on (tangible capital and employment) -> \textcolor{red}{log tangible capital and log employment} \textcolor{orange}{[OK: the code regresses log VA on log employment and log tangible fixed assets, demeaned by year; it also runs the regression separately by 2-digit sector, which the text does not say]} with year fixed effects. More sophisticated total factor productivity measures are feasible but they are highly correlated with simple OLS-based estimates \citep{van2012total} and therefore unlikely to change our results. We compute labor share as the ratio of total labor costs to value added.

Our main analysis measures firm-level productivity and labor share as three-year averages centered on $t^D_i-4$ (i.e.\ over $t^D_i-5$, $t^D_i-4$ and $t^D_i-3$). For each displaced worker this is a fixed pre-displacement window relative to the worker's own event year, so our estimates reflect displacement-induced reallocation of workers to firms with different initial productivity rather than within-firm changes following the layoff. As a robustness exercise, we also recompute firm-level outcomes as within-firm averages over the full 2002--2015 period (Appendix Figures~\ref{fig:prod_et_al_Full_LE} and \ref{fig:prod_et_al_Full_MLO}); the qualitative pattern is preserved and quantitatively amplified.


\subsection{Descriptive statistics on the estimating sample} 
\label{sec:stat_desc}
Our main sample consists of 42,264 individuals equally distributed among treated and control individuals. Standard descriptive statistics on the quantitative variables used in the matching process are reported in Table \ref{tab:desc:AllLE}, where all variables are measured before displacement. Displaced individuals at baseline are 37 years old on average, earn 15.27 euros per hour and worked 1,594 hours. As expected given the matching procedure described above, these variables are well balanced across groups: the Imbens-Rubin normalized difference reported in column (3) never exceeds 0.25, in line with \citet{rubin_assessing_2015} rule of thumb. %suggesting that the matching procedure is successful in reaching balance along matching variables. 

Table \ref{tab:desc:AllLE} also reports the distribution of quantitative variables that were not used in the matching process. About 60\% of the sample are males. They earned about 24,800 euros when considering all their contracts. We also see that the two groups were, on average, employed in similar type of firms in terms of the 2002-2015 wage premia of the pre-layoff employers. %-- estimated from the AKM framework using exhaustive data of firms over the 2001-2015 period
% with godechot, starts in 2002
 In these dimensions as well, both treated and control groups look similar according to the Imbens and Rubin normalized difference. %Here again, none of the differences in distribution seems to exceed Rubin and Imbens rule of thumb. 

Figure \ref{fig:desc:LEquali} finally displays the distribution of LE workers by categories of sector of employment (\ref{fig:desc:LEsec}), occupation (\ref{fig:desc:LEocc}) and highest diploma (\ref{fig:desc:LEdip}). Sector is a matching variable so the displaced and control distributions coincide at $t^D-1$ by construction; we therefore plot the control group only and label it as the \textit{estimation sample}, with the broader DADS panel population shown for reference. Occupation and educational attainment are not used in the matching algorithm and are displayed as a Displaced-vs-Control comparison; both are well balanced. One year before layoff, displaced and control workers mainly worked in manufacturing and extraction or the retail sectors and a large share were blue collars or clerical workers. Compared to the broader DADS panel, our estimation sample is over-represented in manufacturing and construction and under-represented in services to firms and public/education/health, reflecting the sectoral incidence of LE events. %(called ``intermediate professions'' in the French classification) or employees. 
The information about diploma is missing for a large part of the sample (and we exclude it from further analysis) but the balance among observed individual remains (with the highest observed share being vocational basic training). Overall we reach a satisfactory level of balance across both matched and unmatched variables (as is usual in the job displacement literature).  


